\chapter{אבני הבניין של CrewAI}
\section{Agent: עובד דיגיטלי בעל תפקיד}
ב-CrewAI ה-Agent אינו עוד קריאה גולמית למודל שפה, אלא ישות מובנית המגלמת עובד דיגיטלי בעל זהות ברורה. כל Agent מוגדר באמצעות שלושה רכיבים מרכזיים המעצבים את התנהגותו: role, המתאר את התפקיד המקצועי שהוא ממלא, למשל חוקר, מהנדס תוכנה או עורך לשוני; goal, המנסח את המטרה הקונקרטית שאליה הוא חותר; ו-backstory, המספק רקע ונקודת מבט ומעניק לו אישיות, סגנון והקשר התנהגותי. צירוף שלושת הרכיבים אינו קישוט ספרותי בלבד, אלא מנגנון מעשי המעצב את ה-prompt הפנימי, ממקד את היגיון המודל ומפחית סטייה מן המשימה. מעבר להגדרה הטקסטואלית, ל-Agent מוצמדים tools — כלים חיצוניים כגון חיפוש ברשת, קריאת קבצים, הרצת קוד או פנייה ל-API — המרחיבים את יכולותיו אל מעבר לטקסט שהמודל מסוגל לייצר מתוך עצמו. כך ה-Agent הופך לישות הפועלת בעולם, ולא רק מדברת עליו. ההבחנה בין Agent לבין model call פשוטה אך עקרונית: model call הוא אירוע חד-פעמי וחסר זיכרון, ואילו ה-Agent הוא יחידה בעלת תפקיד מתמשך, יכולות מוגדרות ואחריות ברורה בתוך ה-Crew. תפיסה זו מאפשרת לפרק מערכת מורכבת לעובדים מתמחים, שכל אחד מהם אחראי לתחום צר ומוגדר, ובכך לשפר את המודולריות, השקיפות ויכולת התחזוקה של הפתרון כולו.

\section{Task: יחידת עבודה מדידה}
ה-Task הוא יחידת העבודה האטומית של CrewAI, והוא מתרגם כוונה כללית למשימה קונקרטית, מוגדרת ומדידה. כל Task נשען על שני רכיבים שאי אפשר לוותר עליהם: description, המתאר במדויק מה צריך להיעשות ואילו שיקולים יש לקחת בחשבון, ו-expected output, המגדיר במפורש כיצד נראית תוצאה ראויה ומה צורתה הצפויה. ההגדרה המפורשת של ה-expected output היא לב העניין, משום שהיא הופכת את המשימה מעמומה לכזו שניתן לאמת מולה באמצעות validation. כאשר התוצר הצפוי הוא, למשל, טבלה בפורמט מסוים, קובץ LaTeX תקין או רשימת ממצאים מובנית, ה-Crew יכול לבדוק אם הפלט עומד בציפייה, במקום להסתפק בהתרשמות סובייקטיבית. כל Task משויך בדרך כלל ל-Agent מסוים האחראי לביצועו, וכך נוצר חיבור ישיר בין התפקיד לבין העבודה המבוצעת בפועל. ה-Task יכול גם לקבל Context מ-tasks שקדמו לו ולהזין את אלה שיבואו אחריו, וכך הוא הופך לחוליה בשרשרת ולא לאי מבודד. הגדרת העבודה כיחידות מדידות תומכת ישירות ב-observability: כל שלב ניתן למעקב, לרישום ולבחינה, וכשל מקומי ניתן לאיתור ולתיקון בלי לפרק את המערכת כולה. כך CrewAI מאפשר לבנות תהליכים אמינים, שבהם ההצלחה אינה עניין של מזל אלא תוצאה של דרישות ברורות וקריטריונים שניתן למדוד מולם בכל ריצה.

\section{Crew, Process ו-Context: הרכבה ודבק}
לאחר שהוגדרו ה-Agents וה-Tasks, נדרש מנגנון שירכיב אותם למערכת פועלת אחת, וכאן נכנסים לתמונה ה-Crew, ה-Process וה-Context. ה-Crew הוא המעטפת המאגדת קבוצת Agents ואת ה-Tasks המוטלים עליהם ליחידת עבודה אחת בעלת מטרה משותפת, בדומה לצוות אנושי שכל חבר בו תורם את התמחותו. ה-Process מגדיר כיצד הצוות מופעל: ב-Process מסוג sequential ה-Tasks מתבצעים בזה אחר זה לפי סדר קבוע, ואילו ב-Process מסוג hierarchical מתמנה מנהל המחלק ומתאם את העבודה בין החברים. בחירת ה-Process קובעת את אופי התיאום ואת זרימת השליטה בין השלבים. הרכיב המאחד את הכול הוא ה-Context, המשמש דבק בין השלבים. ה-Context הוא המידע העובר משלב לשלב: תוצר של Task אחד הופך לקלט או לרקע של ה-Task הבא, כך שכל שלב נבנה על גבי מה שקדם לו ולא מתחיל מאפס. בלי Context היה הצוות אוסף של עובדים מנותקים שאינם מודעים זה לעבודתו של זה; עם Context נוצרת שרשרת קוהרנטית שבה הידע מצטבר וזורם. כך הרכבת ה-Agents וה-Tasks בתוך Crew, יחד עם Process מתאים ו-Context רציף, היא שמאפשרת ל-CrewAI להפוך רכיבים בודדים ל-orchestration של ממש, שבו השלם גדול מסכום חלקיו ומסוגל לבצע משימות מורכבות מקצה לקצה.

לקריאה נוספת בנושאי פרק זה, ראו \cite{langchain_docs}.

